\chapter{Tier 1: deterministic investment appraisal}\label{chap:tier1}
\section{Problem analysis}\label{sec:problem1}
We select battery and electrolyser capacities and their daily use over
2020--25 (2,192 consecutive days). The investment is made once and applies
to all six years. Direct electricity sales provide the no-investment
benchmark. The investment case depends on three questions: whether
electrolysis can use electricity that could otherwise be sold, how much
hydrogen a buyer can take, and whether storage can shift enough electricity
to recover its costs.

\section{Research findings and model inputs}\label{sec:research1}
\begin{table}[H]
\centering\small
\caption{Research findings used in the model. Appendix~\ref{app:research} gives the evidence and calculations.}
\label{tab:findings}
\begin{tabularx}{\linewidth}{@{}lY@{}}
\toprule
Research subject & Finding used in Tier 1 \\
\midrule
Electricity sales & No client-specific shortfall charge identified. Model uses market receipts. \\
Hydrogen value & UK support benchmark implies a scenario netback of \gbp 7.80/kg \cite{har}. \\
Hydrogen offtake & UK use implies a broad case of 1,559 t/day \cite{eho}. \\
Battery equipment & Four short-duration units with 90--92\% round-trip efficiency \cite{v6}. \\
Electrolyser equipment & Eight units use 55.9--59.0 kWh/kg at rated load \cite{v6}. \\
Equipment economics & Include battery capacity upkeep and electrolyser stack renewals \cite{beis,v6}. \\
Price basis & Convert market prices and equipment costs to 2020 pounds \cite{cpi,basis}. \\
\bottomrule
\end{tabularx}
\end{table}
\begin{table}[H]
\centering\small
\caption{Per-unit equipment inputs. Capital and fixed operation are in 2020 pounds.}
\label{tab:equipment}
{\footnotesize\setlength{\tabcolsep}{2.5pt}
\begin{minipage}[t]{.49\linewidth}
\centering
\textit{(a) Batteries}\par\vspace{3pt}
\begin{tabular}{@{}lrrrrr@{}}
\toprule
Model & Energy & Power & RTE & Capital & FOM \\
 & (MWh) & (MW) & (\%) & (\gbp k) & (\gbp k/yr) \\
\midrule
ST255CS & 0.257 & 0.125 & 90 & 65.6 & 1.90 \\
LUNA215 & 0.215 & 0.108 & 90 & 55.1 & 1.60 \\
LUNA241 & 0.241 & 0.108 & 90 & 60.8 & 1.76 \\
Megapack 2XL & 3.854 & 1.927 & 92 & 1,239 & 31.0 \\
\bottomrule
\end{tabular}
\par\smallskip
\raggedright
RTE is round-trip efficiency. FOM is fixed operation and maintenance.
Use 20--80\% state of charge and 15 years of service. The initial
20\% fill is a commissioning cost. Main-case daily self-discharge is zero.
\end{minipage}\hfill
\begin{minipage}[t]{.49\linewidth}
\centering
\textit{(b) Electrolysers}\par\vspace{3pt}
\begin{tabular}{@{}lrrrr@{}}
\toprule
Model & Input & Output & Consumption & Min. load \\
 & (MW) & (kg/h) & (kWh/kg) & (\%) \\
\midrule
Neptune II & 2.00 & 35 & 57 & 25 \\
Neptune V & 5.03 & 90 & 55.9 & 12.5 \\
Alpha 50 & 50.31 & 900 & 55.9 & 12.5 \\
MC250 & 1.31 & 22.125 & 59 & 10 \\
MC500 & 2.54 & 44.25 & 57.3 & 10 \\
MC1000 & 5.02 & 88.5 & 56.77 & 20 \\
MC2000 & 10.05 & 177 & 56.77 & 20 \\
MC4000 & 20.10 & 354 & 56.77 & 20 \\
\bottomrule
\end{tabular}
\par\smallskip
\raggedright
Use a 30-year life. Per MW, capital is \gbp 1.46m, fixed operation
is \gbp 25.9k/year and each stack renewal costs \gbp 489.6k
in years 11 and 22.
\end{minipage}}
\end{table}


\section{Notation and assumptions}\label{sec:notation1}
\begingroup
\small
\setlength{\tabcolsep}{3pt}
\renewcommand{\arraystretch}{.98}
\begin{tabularx}{\linewidth}{@{}lY@{}}
\toprule
Symbol & Definition \\
\midrule
\(t\in\mathcal T,\ T,\ Y\) & UTC day, number of days, represented years from calendar-year fractions. \\
\(G_t,D_t,S_t,p_t\) & Generation, market demand, surplus \(\max(0,G_t-D_t)\), and price (\gbp/MWh). \\
\(b_i,n_j,z_{jt}\) & Installed battery and electrolyser counts and online electrolyser count. \\
\(q_t,g_t^H,h_{jt}\) & Direct sales, direct input to electrolysis and input to model \(j\). \\
\(c_{it},r_{it},u_{it},s_{it}\) & Charge, discharge to sales, discharge to electrolysis and usable closing stock. \\
\(W_i,P_i,\eta_i,\delta_i\) & Usable capacity, power (MW), one-way efficiency (\(\sqrt{\mathrm{RTE}}\)) and daily self-discharge. \\
\(H_j,\alpha_j,\kappa_j\) & Full-load daily input per unit, minimum load fraction and specific use (MWh/kg). \\
\(v,w_j,L\) & Hydrogen netback and stack reserve (\gbp/kg), and offtake limit (kg/day). \\
\(a_i^B,a_j^H\) & Annual costs per battery and electrolyser unit (\gbp/year). \\
\bottomrule
\end{tabularx}
\endgroup

\noindent Energy and stock are in MWh. Indices \(i,j\) denote battery and electrolyser models.

{\interlinepenalty=10000
\textbf{Assumptions.}
Fixed generation combines the reported total and simulated wind farm.
The single market caps sales at demand, with no imports, other exports,
spatial limits, unused-power revenue or shortfall charge \cite{brief,teacher}.
A dispatches by value. B prioritises sales from generation, then storage,
even at negative prices. Hydrogen is sold daily at \gbp 7.80/kg after
non-electricity variable costs, capped at 1.56, 10 or 1,559 t/day.
Costs use 2020 pounds, a 5\% real discount and zero residual value.\par}

\section{Mathematical model}\label{sec:model1}
\subsection{Objective}
\begin{equation}
\max\quad \Pi =
\frac{1}{Y}\sum_{t\in\mathcal T}
\left[p_t\left(q_t+\sum_i r_{it}\right)+
\sum_j (v-w_j)\frac{h_{jt}}{\kappa_j}\right]
-\sum_i a_i^B b_i-\sum_j a_j^H n_j .
\label{eq:profit}
\end{equation}
Annual cost \(a\) uses capital \(K\), fixed operation \(F\), life \(N\),
discount rate \(d\) and renewals \(C_k\) in year \(k\).
\begin{equation}
a=F+\frac{d}{1-(1+d)^{-N}}
\left(K+\sum_k\frac{C_k}{(1+d)^k}\right).
\label{eq:annualcost}
\end{equation}
Set \(w_j=0\) when renewals enter \(a\). Gain is profit above the
same-policy no-investment benchmark, so existing asset costs cancel.

\subsection{Constraints}
Daily constraints hold for all applicable \(t,i,j\), with \(T=|\mathcal T|\)
and \((x)_+=\max(0,x)\).
{\setlength{\jot}{6pt}
\begin{gather}
 q_t+g_t^H+\sum_i c_{it}\le G_t,\quad
 q_t+\sum_i r_{it}\le D_t,\quad
 \sum_j h_{jt}=g_t^H+\sum_i u_{it}, \label{eq:flow}\\
 s_{it}=(1-\delta_i)s_{i,t-1}+\eta_i c_{it}-(r_{it}+u_{it})/\eta_i,
 \quad 0\le s_{it}\le W_i b_i,\quad s_{i0}=s_{iT}=0, \label{eq:stock}\\
 c_{it}\le b_i\min(24P_i,W_i/\eta_i),\quad
 r_{it}+u_{it}\le b_i\min(24P_i,\eta_i W_i), \label{eq:window}\\
 (c_{it}=0)\ \lor\ (r_{it}+u_{it}=0), \label{eq:direction}\\
 \alpha_j H_j z_{jt}\le h_{jt}\le H_j z_{jt},\quad
 z_{jt}\le n_j,\quad \sum_j h_{jt}/\kappa_j\le L, \label{eq:hydrogen}\\
 \text{Policy B:}\quad q_t=\min(G_t,D_t), \quad
 \sum_i r_{it}=\min\!\left\{(D_t-G_t)_+,\
                   \sum_i\eta_i(1-\delta_i)s_{i,t-1}\right\}. \label{eq:priority}
\end{gather}}
Flows are nonnegative and \(b_i,n_j,z_{jt}\in\mathbb Z_{\ge0}\).

{\setlength{\parskip}{0pt}
\noindent\textbf{Constraint interpretation.}\par
\noindent\eqref{eq:flow} allocates generation, caps sales and balances electrolysis inputs.\par
\noindent\eqref{eq:stock}--\eqref{eq:direction} carry stock across days with losses, limit
throughput and allow one direction per battery model per day.\par
\noindent\eqref{eq:hydrogen} imposes online loading and the daily offtake cap.\par
\noindent\eqref{eq:priority} gives B's remaining demand first claim on usable stock.
The batteries can discharge that stock in one day. Binary indicators
implement the direction choice and B's minimum.\par}

\textbf{Numerical checks.} HiGHS certified 14 cases (relative gap below
\(10^{-9}\)). Independent checks covered flows, cash and dispatch,
including six forced-storage cases \cite{v6}.

\newpage
\section{Results and analysis}\label{sec:results1}
\subsection{Capacity and historical performance}
In the broad market case, electrolysis-only and joint investment return
the same portfolios. A selects 3,631 MW (721 Neptune V and two Neptune II
units) and B selects 740 MW (147 Neptune V units). At 10 t/day, both choose
23.4 MW across six smaller modules. Initial capital in the broad case is
\gbp 5.30bn (A) and \gbp 1.08bn (B), producing about 488,000 and
19,100 t/year of hydrogen respectively.
\begin{table}[H]
\centering\small
\caption{Six-year portfolio choices and annual profit gains by policy.
A dash marks an excluded investment. Gains are in \gbp m/year.}
\label{tab:annual}
\textit{(a) Allowed assets at the broad market cap}\par\vspace{2pt}
\begin{tabular}{@{}l r r r r r r@{}}
\toprule
 & \multicolumn{3}{c}{Policy A} & \multicolumn{3}{c}{Policy B} \\
\cmidrule(lr){2-4}\cmidrule(l){5-7}
Allowed investment & Battery & Electrolysis & Gain & Battery & Electrolysis & Gain \\
 & (MWh) & (MW) & & (MWh) & (MW) & \\
\midrule
Battery only & 0 & -- & 0 & 0 & -- & 0 \\
Electrolysis only & -- & 3,631 & 1,639 & -- & 740 & 37.71 \\
Joint choice & 0 & 3,631 & 1,639 & 0 & 740 & 37.71 \\
\bottomrule
\end{tabular}
\par\vspace{5pt}
\textit{(b) Smaller outlets with both assets allowed}\par\vspace{2pt}
\begin{tabular}{@{}l r r r r r@{}}
\toprule
Daily offtake & Battery & \multicolumn{2}{c}{Electrolysis (MW)} &
 \multicolumn{2}{c}{Gain} \\
\cmidrule(lr){3-4}\cmidrule(l){5-6}
(t/day) & (MWh), A/B & A & B & A & B \\
\midrule
1.56 & 0 & 3.84 & 3.84 & 1.91 & 0.35 \\
10 & 0 & 23.42 & 23.42 & 12.68 & 2.44 \\
\bottomrule
\end{tabular}
\end{table}

At 10 t/day, B loses money in 2021, 2023 and 2024, with a worst
year of \gbp 1.26m below no investment. These years have 43, 29 and
30 surplus days, respectively, against 154 in 2022. A can produce on
additional days when the hydrogen value exceeds
the price of forgone electricity sales. At the broad cap, A meets 82.2\%
of electricity demand and B meets 92.2\%. Using the alternative UK cost
basis at that cap leaves A's capacity unchanged and raises B's from
740 to 1,465 MW, with zero batteries in both cases \cite{beis,v6}.

\subsection{Energy attribution and contract sensitivity}\label{sec:contract1}
The historical mean physical surplus is 2,276 GWh/year.
For the zero-battery solutions we attribute \(\min(g_t^H,S_t)\) to physical
surplus first. The remaining electrolysis input uses energy within current
electricity demand. We calculate the monetary loss separately against
same-policy electricity sales without investment, accounting for negative
prices under A. The aggregate records identify the available fuels by day
and cannot trace the generating unit that supplied electrolysis.
\begin{table}[H]
\centering\small
\caption{Electrolysis sourcing and forgone electricity receipts.
Energy is in GWh/year. The revenue reduction is in \gbp m/year.}
\label{tab:attribution}
\begin{tabular}{@{}l c r r r r@{}}
\toprule
Offtake & Policy & \shortstack{Total input} & \shortstack{Physical\\surplus} &
\shortstack{Within current\\demand} & \shortstack{Electricity\\revenue reduction} \\
\midrule
Broad & A & 27,255 & 2,184 & 25,070 & 1,618 \\
Broad & B & 1,067 & 1,067 & 0 & 0 \\
10 t/day & A & 183 & 43 & 140 & 9.18 \\
10 t/day & B & 43 & 43 & 0 & 0 \\
\bottomrule
\end{tabular}
\end{table}

In the broad market case, A uses about 2,184 GWh/year of physical surplus,
only 8\% of its total electrolysis input. Its electricity receipts fall by
\gbp 1,618m/year. B uses 1,067 GWh/year of surplus and preserves direct
sales, limiting its hydrogen throughput. At 10 t/day, each captures
about 43 GWh/year, just under 2\% of the physical surplus.

\newpage
\noindent
After reoptimising electrolysis and dispatch, the lowest netbacks that
support any investment are about \gbp 3.32/kg in A and \gbp 4.57/kg in B
at 10 t/day. At 1.56 t/day they rise to \gbp 3.38/kg and \gbp 4.66/kg.
The 23.4 MW portfolio in Table~\ref{tab:annual} is sized at \gbp 7.80/kg.
Holding its operation fixed, it breaks even at \gbp 3.90/kg in A or
\gbp 4.61/kg in B. At other prices, reoptimise operation and capacity.

\subsection{Storage and investment sensitivity}\label{sec:storage1}
\begin{table}[H]
\centering\small
\caption{Storage entry under historical and stressed conditions.
Price multipliers scale the standard deviation of daily prices.}
\label{tab:storage}
\begin{tabularx}{\linewidth}{@{}P{.35\linewidth}Y@{}}
\toprule
Test & Implication for storage investment \\
\midrule
Historical high-volatility windows &
All four battery options remain unselected in the highest-variance
60-, 90-, 180- and 365-day windows under both policies. The 90-day price
standard deviation is \gbp 114/MWh against \gbp 65/MWh over six years. \\
Battery cost and price stress &
Halving purchase and maintenance costs still gives zero batteries.
With historical supply and demand, batteries first enter at about
\(40\times\) (A) and \(127\times\) (B) the historical price standard
deviation at the modelled costs. \\
Joint supply and price stress &
The joint grid selects zero batteries at the historical price dispersion
across nine supply--demand gap dispersions. At \(120\times\) price
dispersion in B, raising gap dispersion from \(1\times\) to \(1.1\times\)
changes the optimum from zero to 2.60 GWh. \\
Cavern compressed air &
A separate 50 MW / 1.5 GWh test, at one-third of the reference battery
installation cost per MWh and 62\% round-trip efficiency, needs
\(5.38\)--\(5.50\times\) (A) or \(15.00\)--\(15.13\times\) (B) price
dispersion at the original supply--demand gap. \\
\bottomrule
\end{tabularx}
\end{table}
The volatility tests centre each series before scaling, preserving
means and date order. The battery and cavern tests allow the electrolyser
choice to adjust. For battery-only investment, even free charging and
the full usable discharge on every positive-price day yield at most
57\% of annual ownership cost across the four models. At historical
prices the cavern's annual operating
gain is \gbp 1.38m in A and \gbp 2.35m in B, against an annual cost of
\gbp 8.52m \cite{stress}.

\subsection{Investment recommendations}\label{sec:recommendations1}
\textbf{Size electrolysis to secured demand.}
For 10 t/day, appraise the 23.4 MW portfolio against the \gbp 34.2m
initial cost. A and B yield \gbp 12.68m and \gbp 2.44m/year respectively
(Table~\ref{tab:annual}). Confirm the ability to displace electricity
sales before relying on A. Under B, use the observed loss years to test
cash reserves and the case for staged procurement. Reduce the appraisal to roughly
3.84 MW if the outlet is near 1.56 t/day. At a netback below the
market entry levels above, defer electrolysis at the tested costs.
The broad cap describes market scale. Capacity at that scale warrants
appraisal only with much larger contracted offtake.

\textbf{Defer daily storage.}
Both cost and historical stress tests support zero batteries
(Table~\ref{tab:storage}). Reopen procurement if a lower installed cost
or additional revenue brings the expected operating gain above annual
ownership costs. For the cavern, the historical gain implies an annual
cost ceiling of \gbp 1.38m under A or \gbp 2.35m under B on these inputs.
